% généré depuis source/11_conclusion_generale.md — ne pas modifier à la main
\chapter*{Conclusion générale}
\addcontentsline{toc}{chapter}{Conclusion générale}
\frontmark{Conclusion générale}

Cette première version de l'étude a appliqué l'ensemble de la démarche de sûreté de fonctionnement à une flotte de turboréacteurs simulés, des données brutes jusqu'à un outil d'aide à la décision.

L'analyse de fiabilité, menée en tenant compte des 100 moteurs censurés, montre une flotte soumise à l'usure, avec une période sans défaillance. La loi de Weibull à trois paramètres, validée par le test de Kolmogorov-Smirnov et équivalente à la loi log-normale selon l'AIC, donne un MTTF de 206,3 cycles et des durées $B_{10}$ et $B_{50}$ de 151,9 et 199,9 cycles. Sous des hypothèses de réparation et de coût explicitement signalées, la disponibilité intrinsèque atteint 95,98 \% et un remplacement préventif à 131 cycles réduit le coût par cycle de 68 \%.

La chaîne de maintenance prédictive complète cette vision globale par une décision moteur par moteur. Un indicateur de santé construit sur 12 capteurs décrit la dégradation de façon monotone, le classifieur détecte l'état critique avec une AUC de 0,991, et deux modèles estiment la durée de vie résiduelle. Le LSTM est le plus précis (RMSE de 15,35 cycles). Le modèle de dégradation bayésien, moins précis (RMSE de 24,47 cycles), est interprétable et probabiliste ; son biais et la calibration insuffisante de ses intervalles sont identifiés, mesurés et expliqués.

Enfin, l'ensemble est livré sous la forme d'un système reproductible : pipeline, API, tableau de bord et suite de tests automatiques, qui garantit que chaque chiffre de ce rapport provient du code. Les perspectives prioritaires pour la Version 2 sont l'amélioration du modèle de dégradation et de ses intervalles, l'estimation de la loi à trois paramètres avec censure, et l'extension aux sous-ensembles comportant plusieurs conditions opératoires.
